\documentclass[11pt,letterpaper]{article}
\def\SuiteDocTitle{Chapter 20: The Release Plan}
\def\SuiteDocCode{SG-20}
% Shared professional style: pdfLaTeX compatible, no external font files.
\usepackage[T1]{fontenc}
\usepackage[utf8]{inputenc}
\usepackage{lmodern}
\usepackage{microtype}
\usepackage[margin=0.83in,headheight=16pt,footskip=28pt]{geometry}
\usepackage[table]{xcolor}
\definecolor{navy}{HTML}{12304A}
\definecolor{teal}{HTML}{087F83}
\definecolor{gold}{HTML}{B78633}
\definecolor{blueback}{HTML}{F0F5F8}
\definecolor{tealback}{HTML}{EEF8F6}
\definecolor{amberback}{HTML}{FFF7E8}
\definecolor{grayback}{HTML}{F6F7F8}
\usepackage{array,tabularx,longtable,booktabs}
\usepackage{enumitem}
\usepackage{amssymb}
\usepackage{titlesec}
\usepackage{fancyhdr,lastpage}
\usepackage{tcolorbox}
\tcbuselibrary{breakable}
\usepackage{pgffor}
\usepackage{hyperref}
\hypersetup{colorlinks=true,linkcolor=teal,urlcolor=teal,pdfauthor={Engineering Learning Series},pdftitle={\SuiteDocTitle}}
\setlength{\parindent}{0pt}
\setlength{\parskip}{5.5pt}
\setlist{topsep=3pt}
\setcounter{tocdepth}{2}
\titleformat{\section}{\Large\bfseries\color{navy}}{\thesection}{0.7em}{}
\titleformat{\subsection}{\large\bfseries\color{teal}}{\thesubsection}{0.6em}{}
\titleformat{\subsubsection}{\normalsize\bfseries\color{navy}}{\thesubsubsection}{0.6em}{}
\pagestyle{fancy}
\fancyhf{}
\fancyhead[L]{\scriptsize\color{navy}AGILE REQUIREMENTS ENGINEERING}
\fancyhead[R]{\scriptsize\color{teal}\SuiteDocCode}
\fancyfoot[L]{\scriptsize Independent learning and implementation materials}
\fancyfoot[R]{\scriptsize Page \thepage\ of \pageref{LastPage}}
\renewcommand{\headrulewidth}{0.35pt}
\renewcommand{\footrulewidth}{0pt}
\newcommand{\SuiteHeader}[2]{%
{\color{teal}\small\bfseries ENGINEERING DOCUMENTATION SUITE \hfill #2}\par\vspace{8pt}
{\color{navy}\LARGE\bfseries\SuiteDocTitle}\par\vspace{5pt}
{\color{black!68}\normalsize #1}\par\vspace{8pt}
{\color{gold}\rule{\linewidth}{2pt}}\par\vspace{8pt}}
\newcommand{\WritingLines}[1]{\par\foreach \n in {1,...,#1}{\noindent\rule{\linewidth}{0.35pt}\par\vspace{0.28cm}}}
\newcommand{\FreeSpace}[1]{\begin{tcolorbox}[colback=white,colframe=black!25,height=#1,boxrule=0.5pt]\end{tcolorbox}}
\newcommand{\Field}[1]{\noindent\textbf{#1}\quad\rule{0.42\linewidth}{0.35pt}\par\vspace{0.3cm}}
\newcommand{\CheckItem}[1]{\noindent$\square$\quad #1\par\vspace{2pt}}
\newcommand{\SmallNote}[1]{{\footnotesize\color{black!62}#1}\par}
\newcommand{\Guidance}[1]{\begin{tcolorbox}[colback=blueback,colframe=teal,breakable,title=Guidance]#1\end{tcolorbox}}
\sloppy

\begin{document}
\SuiteHeader{An Example | Book pp. 219–222}{Study guide}
\tableofcontents\medskip\noindent\rule{\linewidth}{0.4pt}\medskip
\section{Learning outcomes}
\begin{itemize}[leftmargin=*,itemsep=2pt]
\item Explain the chapter's principal argument without confusing a story with a finished specification.
\item Apply the chapter's guidance to a traceable, user-valued increment.
\item Identify a failure mode, select evidence for a correction, and defend the result.
\end{itemize}
\section{Chapter synthesis}
\subsection{Estimate capacity and select outcomes}
The book synthesizes story estimates and value priorities into a proposed release.

\begin{tcolorbox}[colback=tealback,colframe=navy,title={Engineering translation},breakable]
Build an increment that allows submission, review and status lookup before secondary enhancements.

\end{tcolorbox}
\begin{tcolorbox}[colback=amberback,colframe=navy,title={Failure mode to watch},breakable]
A release full of technical tasks may still lack usable workflows.

\end{tcolorbox}
\subsection{Manage risk visibly}
High-value and high-risk work both affect sequencing.

\begin{tcolorbox}[colback=tealback,colframe=navy,title={Engineering translation},breakable]
Place the unknown identity integration experiment before the dependent authentication slice.

\end{tcolorbox}
\begin{tcolorbox}[colback=amberback,colframe=navy,title={Failure mode to watch},breakable]
Do not overload the first iteration with unbounded discovery.

\end{tcolorbox}
\subsection{Adapt as evidence changes}
The plan is a negotiated forecast, not a permanent contract.

\begin{tcolorbox}[colback=tealback,colframe=navy,title={Engineering translation},breakable]
Review realized velocity and acceptance outcomes at each iteration boundary.

\end{tcolorbox}
\begin{tcolorbox}[colback=amberback,colframe=navy,title={Failure mode to watch},breakable]
Keep a change log explaining material replanning.

\end{tcolorbox}
\section{Applied case: Service Request Portal}
The running example is original to this suite, not an example reproduced from the book. Its users submit, review, and audit requests. The proposed story below is a starting point for discussion, not a complete contract.

\begin{tcolorbox}[colback=blueback,colframe=navy,title={Candidate story},breakable]
As a requester, I want to know when my request changes state so that I can respond promptly.

\end{tcolorbox}
\section{Structured workshop}
\begin{enumerate}[leftmargin=*,itemsep=3pt]
\item Read the candidate story aloud and ask the named role what success looks like.
\item Capture at least three decisions and at least two unanswered questions.
\item Identify a valid path, one exception, and an observable result; convert each relevant distinction into acceptance evidence.
\item Check whether the result is independently demonstrable and appropriately sized for a near-term iteration.
\item Record the responsible decision owner and the next date or event for validation.
\end{enumerate}
\section{Practice exercises}
\begin{enumerate}[leftmargin=*,itemsep=3pt]
\item Produce a three-iteration release scenario under a point budget.
\item Describe the scope cut if one integration takes longer than expected.
\end{enumerate}
\section{Knowledge checks}
\begin{enumerate}[leftmargin=*,itemsep=3pt]
\item What evidence shows that the work meets a real user goal rather than just a technical objective?
\item Which assumption could most invalidate the story, and how would you test that assumption early?
\item What artifact or conversation would you preserve for traceability, and what could be safely discarded?
\item How would the approach differ for a near-term story versus a distant idea?
\end{enumerate}
\section{Completion evidence}
\begin{itemize}[leftmargin=*,itemsep=2pt]
\item One revised story with named role, goal, value, owner and open questions.
\item A concrete acceptance example that could fail on incorrect behavior.
\item A short justification for sizing, sequencing or documentation depth.
\item A dated note of feedback received and the revision it caused.
\end{itemize}
\section{Reading reference and scope}
Primary source: Mike Cohn, User Stories Applied: For Agile Software Development (2004), Chapter 20, "The Release Plan," printed pages 219–222. This is an original interpretive study guide, not a chapter transcription. Product examples, workflow controls and traceability suggestions in this guide are instructional extensions.


\end{document}
